% GENERATED FILE. Edit canonical lessons, then run npm run generate:books.
% canonical-source-hash: fnv1a64:43a2be03e46636cc
% canonical-lessons: ES-C546-bautizo, ES-C546-padrino, ES-C546-madrina, ES-C546-nieto, ES-C546-panal

\chapter{El bautizo}
\label{ch:es-el-bautizo}
\hlchaptermodality{546}

\begin{chapteropening}
\textbf{By the end of this chapter:} \emph{I can talk about a christening: the godparents, the baby's nappies and the grandchildren.}

Complete the last lesson of chapter 546: I can talk about a christening: the godparents, the baby's nappies and the grandchildren.
\end{chapteropening}

\section[el bautizo]{el bautizo — the christening}

\label{lesson:ES-C546-bautizo}



\begin{quote}
Before the new one: say the Spanish for adolescence, then the Spanish for the generation.
\end{quote}

\subsection*{You'll want to know: el bautizo}
\textbf{el bautizo} — \textquotedblleft{}the christening\textquotedblright{}.

\begin{quote}
El bautizo fue en la iglesia del pueblo.
\end{quote}

\begin{quote}
\emph{The christening was in the village church.}
\end{quote}

\emph{Bautizar} is to christen. \emph{El bautismo} is the sacrament itself.

\begin{cousinweb}
\emph{bautizo} is built on \emph{bautizar}, which comes from Church Latin \emph{baptīzāre} 'to baptize' (Lewis \& Short), taken from Greek \emph{baptízein} 'to dip in water'.
\end{cousinweb}

\begin{grammarlens}[title={What you've built}]
One more word for telling, explaining and describing at B1.
\end{grammarlens}

\subsection*{Your turn}
\begin{itemize}
\raggedright
  \item \emph{Say it:} \emph{el bautizo}
  \item \emph{Say it:} \emph{el bautizo}, once more
  \item \emph{Say it:} a sentence of your own with \emph{el bautizo}
  \item \emph{Recall:} say the Spanish for adolescence, then the Spanish for the generation, then say \emph{el bautizo} again
\end{itemize}

\begin{tcolorbox}[breakable,colback=teal!4,colframe=teal!35!black,title={Before you move on}]
What does el bautizo mean? (\textquotedblleft{}The christening\textquotedblright{}.) Say it once more.
\end{tcolorbox}
\section[el padrino]{el padrino — the godfather; the best man}

\label{lesson:ES-C546-padrino}



\begin{quote}
Before the new one: say the Spanish for the generation, then the Spanish for the christening.
\end{quote}

\subsection*{You'll want to know: el padrino}
\textbf{el padrino} — \textquotedblleft{}the godfather; the best man\textquotedblright{}.

\begin{quote}
Mi tío es mi padrino.
\end{quote}

\begin{quote}
\emph{My uncle is my godfather.}
\end{quote}

At a wedding, \emph{el padrino} walks with the bride or the groom. \emph{Los padrinos} — the godparents.

\begin{cousinweb}
\emph{padrino} is inherited from spoken Latin \emph{*patrīnus} 'a godfather' (EtymDB), built on \emph{pater} 'father'.
\end{cousinweb}

\begin{grammarlens}[title={What you've built}]
One more word for telling, explaining and describing at B1.
\end{grammarlens}

\subsection*{Your turn}
\begin{itemize}
\raggedright
  \item \emph{Say it:} \emph{el padrino}
  \item \emph{Say it:} \emph{el padrino}, once more
  \item \emph{Say it:} a sentence of your own with \emph{el padrino}
  \item \emph{Recall:} say the Spanish for the generation, then the Spanish for the christening, then say \emph{el padrino} again
\end{itemize}

\begin{tcolorbox}[breakable,colback=teal!4,colframe=teal!35!black,title={Before you move on}]
What does el padrino mean? (\textquotedblleft{}The godfather; the best man\textquotedblright{}.) Say it once more.
\end{tcolorbox}
\section[la madrina]{la madrina — the godmother}

\label{lesson:ES-C546-madrina}



\begin{quote}
Before the new one: say the Spanish for the christening, then the Spanish for the godfather.
\end{quote}

\subsection*{You'll want to know: la madrina}
\textbf{la madrina} — \textquotedblleft{}the godmother\textquotedblright{}.

\begin{quote}
Mi madrina me regaló una bicicleta.
\end{quote}

\begin{quote}
\emph{My godmother gave me a bicycle.}
\end{quote}

At a wedding she is the woman who accompanies the groom.

\begin{cousinweb}
\emph{madrina} is inherited from spoken Latin \emph{*mātrīna} (EtymDB), built on \emph{māter} 'mother'.
\end{cousinweb}

\begin{grammarlens}[title={What you've built}]
One more word for telling, explaining and describing at B1.
\end{grammarlens}

\subsection*{Your turn}
\begin{itemize}
\raggedright
  \item \emph{Say it:} \emph{la madrina}
  \item \emph{Say it:} \emph{la madrina}, once more
  \item \emph{Say it:} a sentence of your own with \emph{la madrina}
  \item \emph{Recall:} say the Spanish for the christening, then the Spanish for the godfather, then say \emph{la madrina} again
\end{itemize}

\begin{tcolorbox}[breakable,colback=teal!4,colframe=teal!35!black,title={Before you move on}]
What does la madrina mean? (\textquotedblleft{}The godmother\textquotedblright{}.) Say it once more.
\end{tcolorbox}
\section[el nieto]{el nieto — the grandson, the grandchild}

\label{lesson:ES-C546-nieto}



\begin{quote}
Before the new one: say the Spanish for the godfather, then the Spanish for the godmother.
\end{quote}

\subsection*{You'll want to know: el nieto}
\textbf{el nieto} — \textquotedblleft{}the grandson, the grandchild\textquotedblright{}.

\begin{quote}
Mis abuelos tienen seis nietos.
\end{quote}

\begin{quote}
\emph{My grandparents have six grandchildren.}
\end{quote}

\emph{La nieta} is the granddaughter; \emph{los nietos} can mean all the grandchildren.

\begin{cousinweb}
\emph{nieto} is usually traced to spoken Latin \emph{*neptus}, a masculine made from Latin \emph{neptis} 'a granddaughter' (Lewis \& Short), which belongs with \emph{nepōs} 'grandson'.
\end{cousinweb}

\begin{grammarlens}[title={What you've built}]
One more word for telling, explaining and describing at B1.
\end{grammarlens}

\subsection*{Your turn}
\begin{itemize}
\raggedright
  \item \emph{Say it:} \emph{el nieto}
  \item \emph{Say it:} \emph{el nieto}, once more
  \item \emph{Say it:} a sentence of your own with \emph{el nieto}
  \item \emph{Recall:} say the Spanish for the godfather, then the Spanish for the godmother, then say \emph{el nieto} again
\end{itemize}

\begin{tcolorbox}[breakable,colback=teal!4,colframe=teal!35!black,title={Before you move on}]
What does el nieto mean? (\textquotedblleft{}The grandson, the grandchild\textquotedblright{}.) Say it once more.
\end{tcolorbox}
\section[el pañal]{el pañal — the nappy, the diaper}

\label{lesson:ES-C546-panal}



\begin{quote}
Before the new one: say the Spanish for the godmother, then the Spanish for the grandson.
\end{quote}

\subsection*{You'll want to know: el pañal}
\textbf{el pañal} — \textquotedblleft{}the nappy, the diaper\textquotedblright{}.

\begin{quote}
Hay que cambiarle el pañal al bebé.
\end{quote}

\begin{quote}
\emph{The baby's nappy needs changing.}
\end{quote}

\emph{Estar en pañales} — to be in its infancy, to know very little about something.

\begin{cousinweb}
\emph{pañal} is built on \emph{paño} 'cloth' (EtymDB), which is inherited from Latin \emph{pannus} 'a piece of cloth'.
\end{cousinweb}

\begin{grammarlens}[title={What you've built}]
One more word for telling, explaining and describing at B1.
\end{grammarlens}

\subsection*{Your turn}
\begin{itemize}
\raggedright
  \item \emph{Say it:} \emph{el pañal}
  \item \emph{Say it:} \emph{el pañal}, once more
  \item \emph{Say it:} a sentence of your own with \emph{el pañal}
  \item \emph{Recall:} say the Spanish for the godmother, then the Spanish for the grandson, then say \emph{el pañal} again
\end{itemize}

\begin{tcolorbox}[breakable,colback=teal!4,colframe=teal!35!black,title={Before you move on}]
What does el pañal mean? (\textquotedblleft{}The nappy, the diaper\textquotedblright{}.) Say it once more.
\end{tcolorbox}
